\section{Годограф Найквиста }

В соответствии с вариантом 27 из \textit{Таблицы 1 инструкции по выполнению л/р} придумаем три объекта пятого
порядка $p$ полюсов передаточных функций которых вещественные, а $q$ - комплексно-
сопряженные.
\begin{itemize}
    \item Первая передаточная функция должна иметь $n$ неустойчивых полюсов у разомкнутой системы и $m$ неустойчивых полюсов у замкнутой.
    \item Вторая передаточная функция должна иметь $0$ неустойчивых полюсов у разомкнутой системы и $m$ у замкнутой.
    \item Третья передаточная функция должна иметь $n$ неустойчивых полюсов у разомкнутой системы и $0$ у замкнутой.
\end{itemize}
Для каждого объекта:
\begin{enumerate}
    \item Приведем карты (изображения размещения на комплексной координатной плоскости) нулей и полюсов разомкнутой и замкнутой систем на комплексной координатной плоскости.
    \item Выполним моделирование и приведем переходные характеристики для разомкну-
    той и замкнутой систем.
    \item Построим годограф Найквиста (АФЧХ). Найдем число оборотов годографа по ча-
    совой стрелке вокруг точки $(-1; 0)$ и через критерий Найквиста. Сравним результаты.
    \item Построим ЛАЧХ и ЛФЧХ разомкнутой системы. Определим суммарное количество переходов ЛФЧХ через критические отрезки при положительном значении
    ЛАЧХ, сравним с результатами (количеством оборотов), полученными в предыдущем пункте. Проанализируем связь устойчивости с видом ЛАФЧХ разомкнутых систем используя логарифмический критерий Найквиста.
\end{enumerate}

\begin{table}[h]
\centering
\begin{tabular}{|c|c|c|c|c|c|c|}
\hline
\textbf{Вариант} & \textbf{p} & \textbf{q} & \textbf{n} & \textbf{m} & \textbf{i} & \textbf{j} \\ \hline
27 & 1 & 4 & 4 & 2 & 13 & 3 \\ \hline
\end{tabular}
\end{table}

Придумаем 3 объекта:
Разомкнутая система имеет вид: $W(s)=\frac{R(s)}{Q(s)}$,
замкнутая - $W_3(s)=\frac{W(s)}{1+W(s)}=\frac{R(s)}{R(s)+Q(s)}=\frac{R(s)}{D(s)}$. Составим с учётом этого системы с полиномами, имеющими подходящие нам корни, найдя $R(s)$ через $Q(s)$ и $D(s)$.

\begin{itemize}
    \item 1:
Математическая модель: $4$ неустойчивых полюса у разомкнутой системы и $2$ неустойчивых полюса у замкнутой.
$$
\begin{cases}
Q(s)= (s+1)(s-1-2i)(s-1+2i)(s-2-3i)(s-2+3i),\\
Q(s)+R(s)= (s+2)(s+1-2i)(s+1+2i)(s-1-3i)(s-1+3i)
\end{cases}\Rightarrow
$$$$
\Rightarrow R(s) = 7 s^4 - 9 s^3 + 52 s^2 + 51 s + 35
$$$$
W(s)=\frac{7 s^4 - 9 s^3 + 52 s^2 + 51 s + 35}{s^5 - 5 s^4 + 20 s^3 - 20 s^2 + 19 s + 65}
$$$$
W_3(s)=\frac{7 s^4 - 9 s^3 + 52 s^2 + 51 s + 35}{s^5 + 2 s^4 + 11 s^3 + 32 s^2 + 70 s + 100}
$$


\item 2:
Математическая модель:
Математическая модель: $0$ неустойчивых полюсов у разомкнутой системы и $2$ неустойчивых полюса у замкнутой.
$$
\begin{cases}
    Q(s)=(s+1)(s+1-i)(s+1+i)(s+2-i)(s+2+i) \\
    Q(s)+R(s)=(s+2)(s+2-3i)(s+2+3i)(s-4+i)(s-4-i)
\end{cases}
$$$$
\Rightarrow R(s) = -9 s^4 - 31 s^3 - 73 s^2 + 121 s + 432
$$$$
W(s)=\frac{-9 s^4 - 31 s^3 - 73 s^2 + 121 s + 432}{s^5 + 7 s^4 + 21 s^3 + 33 s^2 + 28 s + 10}
$$$$
W_3(s)=\frac{-9 s^4 - 31 s^3 - 73 s^2 + 121 s + 432}{s^5 - 2 s^4 - 10 s^3 - 40 s^2 + 149 s + 442}
$$


\item 3:
Математическая модель:
Математическая модель: $4$ неустойчивых полюса у разомкнутой системы и $0$ неустойчивых полюсов у замкнутой.
$$
\begin{cases}
    Q(s)=(s+1)(s-1-i)(s-1+i)(s-2-i)(s-2+i) \\
    Q(s)+R(s)=(s+2)(s+2-4i)(s+2+4i)(s+3-i)(s+3+i)
\end{cases}
$$$$
\Rightarrow R(s) = 17 s^4 + 65 s^3 + 271 s^2 + 528 s + 390
$$$$
W(s)=\frac{17 s^4 + 65 s^3 + 271 s^2 + 528 s + 390}{s^5 - 5 s^4 + 9 s^3 - 3 s^2 - 8 s + 10}
$$$$
W_3(s)=\frac{17 s^4 + 65 s^3 + 271 s^2 + 528 s + 390}{s^5 + 12 s^4 + 74 s^3 + 268 s^2 + 520 s + 400}
$$


\end{itemize}


\textbf{Построим карты для разомкнутой и замкнутой систем (см. рис. \ref{c1}-\ref{c3}).}

\begin{figure}[H]
    \centering
    \includegraphics[width=1\linewidth]{lab_6/images/pz_combined1.png}
    \caption{Карта (объект 1)}
    \label{c1}
\end{figure}
\begin{figure}[H]
    \centering
    \includegraphics[width=1\linewidth]{lab_6/images/pz_combined2.png}
    \caption{Карта (объект 2)}
    \label{c2}
\end{figure}
\begin{figure}[H]
    \centering
    \includegraphics[width=1\linewidth]{lab_6/images/pz_combined3.png}
    \caption{Карта (объект 3)}
    \label{c3}
\end{figure}

Проведем моделирование.\\

\textbf{Переходная характеристика для разомкнутой и замкнутой систем (см. рис. \ref{s1}-\ref{s3}).}
\begin{figure}[H]
    \centering
    \includegraphics[width=0.62\linewidth]{lab_6/images/step_combined1.png}
    \caption{Переходная характеристика (объект 1)}
    \label{s1}
\end{figure}
\begin{figure}[H]
    \centering
    \includegraphics[width=0.62\linewidth]{lab_6/images/step_combined2.png}
    \caption{Переходная характеристика (объект 2)}
    \label{s2}
\end{figure}
\begin{figure}[H]
    \centering
    \includegraphics[width=0.62\linewidth]{lab_6/images/step_combined3.png}
    \caption{Переходная характеристика (объект 3)}
    \label{s3}
\end{figure}

\textbf{Годограф Найквиста (см. рис. \ref{n1}-\ref{n3}).}
\begin{figure}[H]
    \centering
    \includegraphics[width=0.6\linewidth]{lab_6/images/nyquist1.png}
    \caption{Годограф Найквиста (объект 1)}
    \label{n1}
\end{figure}
Число оборотов вокруг точки (против ч.с.) (-1, 0): 2\\
Критерий Найквиста:$m = n + N$,
где\\
1. $m$ –число неустойчивых полюсов замкнутой сисемы\\
2. $n$ – Число неустойчивых полюсов разомкнутой системы\\
3. $N$ –Число оборотов годографа по часовой стрелке вокруг точки (-1,0)\\
Сравнение с результатами по критерию Найквиста: $2=4-2$ - сходится.

\begin{figure}[H]
    \centering
    \includegraphics[width=0.6\linewidth]{lab_6/images/nyquist2.png}
    \caption{Годограф Найквиста (объект 2)}
    \label{n2}
\end{figure}
Число оборотов вокруг точки (против ч.с.) (-1, 0): -2\\
Сравнение с результатами по критерию Найквиста: $2=0+2$ - сходится.

\begin{figure}[H]
    \centering
    \includegraphics[width=0.6\linewidth]{lab_6/images/nyquist3.png}
    \caption{Годограф Найквиста (объект 3)}
    \label{n3}
\end{figure}
Число оборотов вокруг точки (против ч.с.)
(-1, 0): 4\\
Сравнение с результатами по критерию Найквиста: $0=4-4$ - сходится.


\textbf{ЛАЧХ и ЛФЧХ разомкнутой системы (см. рис. \ref{l1}-\ref{l3}).}

\begin{figure}[H]
    \centering
    \includegraphics[width=0.9\linewidth]{lab_6/images/bode1.png}
    \caption{ЛАЧХ и ЛФЧХ (объект 1)}
    \label{l1}
\end{figure}
Суммарное количество переходов ЛФЧХ через критические отрезки при положительном значении ЛАЧХ: 1\\
\textbf{Анализ связи устойчивости с видом ЛАФЧХ разомкнутых систем используя логарифмический критерий Найквиста:}
Для разомкнутой системы число правых корней характеристического уравнения
$r = 4$, следовательно, согласно логарифмическому критерию Найквиста должно
выполняться условие $N_+ - N_- = r/2 = 2$. По графику ЛФЧХ данное условие не выполняется, что подтверждает неустойчивость замкнутой системы.

\begin{figure}[H]
    \centering
    \includegraphics[width=0.9\linewidth]{lab_6/images/bode2.png}
    \caption{ЛАЧХ и ЛФЧХ (объект 2)}
    \label{l2}
\end{figure}
Суммарное количество переходов ЛФЧХ через критические отрезки при положительном значении ЛАЧХ: -1\\
\textbf{Анализ связи устойчивости с видом ЛАФЧХ разомкнутых систем используя логарифмический критерий Найквиста:}
Для разомкнутой системы число правых корней
характеристического уравнения равно $r = 0$, следовательно, согласно логарифмическому критерию Найквиста должно
выполняться условие $N_+ - N_- = 0$. По ЛФЧХ при положительном значении ЛАЧХ наблюдается $-1/2$ переходов через критические отрезки, что подтверждает наличие двух неустойчивых полюсов у замкнутой системы.

\begin{figure}[H]
    \centering
    \includegraphics[width=0.9\linewidth]{lab_6/images/bode3.png}
    \caption{ЛАЧХ и ЛФЧХ (объект 3)}
    \label{l3}
\end{figure}
Суммарное количество переходов ЛФЧХ через критические отрезки при положительном значении ЛАЧХ: 2\\
\textbf{Анализ связи устойчивости с видом ЛАФЧХ разомкнутых систем используя логарифмический критерий Найквиста:}
Для разомкнутой системы число правых корней равно $r = 4$, следовательно, должно выполняться равенство $N_+ - N_- = r/2 = 2$. Это ранво числу переходов ЛФЧХ через критические отрезки и подтверждает устойчивость замкнутой системы.

\textbf{Сравнение с предыдущим пунктом:} для устойчивости замкнутой системы число оборотов годографа Найквиста против часовой стрелки должно совпадать с числом правых корней разомкнутой системы.

\section{Коэффициент усиления}

В соответствии с вариантом 27 из \textit{Таблицы 1 инструкции по выполнению л/р} возьмем значение $i$ и соответствующие ему передаточные функции $W_1(s)$ и $W_2(s)$ (Таблица 2). Добавим к каждой
функции коэффициент усиления $k$. Считаем $k$ положительным. Для полученных систем:
\begin{enumerate}
    \item Построим годограф Найквиста для значения коэффициента усиления $k = 1$.
    \item Рассмотрим, как влияет коэффициент усиления $k$ на кривую годографа, взяв не
    менее трех значений $k$.
    \item Найдем зависимость количества неустойчивых полюсов замкнутой системы относительно значений коэффициента усиления $k$. Найдем пределы значений коэффициента усиления $k$ относительно которых замкнутая система устойчива.
    \item Определим значение запаса устойчивости по амплитуде (если применимо).
    \item Выполним моделирование и приведем переходные характеристики замкнутой системы при значениях коэффициента $k$ для каждой из определенных областей значений, соответствующих как устойчивой системе, так и неустойчивой.
\end{enumerate}

\begin{table}[h]
\centering
\begin{tabular}{|c|c|c|}
\hline
\textbf{i} & $\mathbf{W_1(s)}$ & $\mathbf{W_2(s)}$ \\ \hline
13 & $\displaystyle \frac{s-1}{s^2+3s+1}$ & $\displaystyle \frac{10s^3-3s^2+13s-2}{10s^3+8s^2-5s+4}$ \\ \hline
\end{tabular}
\end{table}





\section{Запаздывание}

В соответствии с вариантом 27 из \textit{Таблицы 1 инструкции по выполнению л/р} возьмем значение $j$ и соответствующие ему передаточные функции $W_3(s)$ и $W_4(s)$ (таблица 3). Добавим к каждой
функции звено чистого запаздывания $e^{-\tau s}$. Для полученных систем:
\begin{enumerate}
    \item Построим годограф Найквиста для значений коэффициента запаздывания $\tau = 0$
    и $\tau = 0.5$.
    \item Рассмотрим, как влияет коэффициент запаздывания $\tau$ на кривую годографа, взяв
    не менее трех значений $\tau$.
    \item Найдем границы значений коэффициента запаздывания $\tau$ относительно которых
    замкнутая система устойчива.
    \item Определим значение запаса устойчивости по фазе (если применимо).
    \item Выполним моделирование и приведем переходные характеристики замкнутой си-
    стемы при значениях коэффициента $\tau$ для каждой из определенных областей значений, соответствующих как устойчивой системе, так и неустойчивой.
\end{enumerate}

\begin{table}[h]
\centering
\begin{tabular}{|c|c|c|}
\hline
\textbf{j} & $\mathbf{W_3(s)}$ & $\mathbf{W_4(s)}$ \\ \hline
3 & $\displaystyle \frac{2s+9}{s^2+s+9}$ & $\displaystyle \frac{10s^2-10s+13}{10s^3+37s+25}$ \\ \hline
\end{tabular}
\end{table}


\newpage
\section{Приложение}

\matlabtitle{Задание 1}
\begin{matlabcode}
% Система 1
R1 = [7 -9 52 51 35];
Q1 = [1 -5 20 -20 19 65];
D1 = [1 2 11 32 70 100];   
open_sys1 = tf(R1, Q1);
closed_sys1 = tf(R1, D1);
p1 = pole(open_sys1);
z1 = zero(open_sys1);
p_closed1 = pole(closed_sys1);
z_closed1 = zero(closed_sys1);

figure('Position', [100 100 1200 500]);
set(gcf, 'Color', 'w');

subplot(1,2,1);
hold on;
grid on;
xlabel('Re');
ylabel('Im');
title('Разомкнутая система 1', 'FontWeight', 'normal');

for i = 1:length(p1)
    if real(p1(i)) < 0
        plot(real(p1(i)), imag(p1(i)), 'xg', 'MarkerSize', 10, 'LineWidth', 2);
    else
        plot(real(p1(i)), imag(p1(i)), 'xr', 'MarkerSize', 10, 'LineWidth', 2);
    end
end
for i = 1:length(z1)
    plot(real(z1(i)), imag(z1(i)), 'bo', 'MarkerSize', 8, 'LineWidth', 2);
end
xlim([-2 3]);
ylim([-4 4]);
box on;
h1 = plot(NaN, NaN, 'xg', 'MarkerSize', 10, 'LineWidth', 2);
h2 = plot(NaN, NaN, 'xr', 'MarkerSize', 10, 'LineWidth', 2);
h3 = plot(NaN, NaN, 'bo', 'MarkerSize', 8, 'LineWidth', 2);
legend([h1, h2, h3], {'Уст. полюс', 'Неуст. полюс', 'Ноль'}, 'Location', 'best');

subplot(1,2,2);
hold on;
grid on;
xlabel('Re');
ylabel('Im');
title('Замкнутая система 1', 'FontWeight', 'normal');

for i = 1:length(p_closed1)
    if real(p_closed1(i)) < 0
        plot(real(p_closed1(i)), imag(p_closed1(i)), 'xg', 'MarkerSize', 10, 'LineWidth', 2);
    else
        plot(real(p_closed1(i)), imag(p_closed1(i)), 'xr', 'MarkerSize', 10, 'LineWidth', 2);
    end
end
for i = 1:length(z_closed1)
    plot(real(z_closed1(i)), imag(z_closed1(i)), 'bo', 'MarkerSize', 8, 'LineWidth', 2);
end
xlim([-3 2]);
ylim([-4 4]);
box on;
h1 = plot(NaN, NaN, 'xg', 'MarkerSize', 10, 'LineWidth', 2);
h2 = plot(NaN, NaN, 'xr', 'MarkerSize', 10, 'LineWidth', 2);
h3 = plot(NaN, NaN, 'bo', 'MarkerSize', 8, 'LineWidth', 2);
legend([h1, h2, h3], {'Уст. полюс', 'Неуст. полюс', 'Ноль'}, 'Location', 'best');

exportgraphics(gcf, "images/pz_combined1.png", Resolution=200);

figure('Position', [100 100 800 800]);
set(gcf, 'Color', 'w');

simtime = 15;
subplot(2,1,1);
[y_open1, t1] = step(open_sys1, simtime);
plot(t1, y_open1, 'b', 'LineWidth', 1.5);
xlabel('t, s');
ylabel('y');
grid on;
title('Разомкнутая система 1', 'FontWeight', 'normal');
xlim([0 simtime]);

subplot(2,1,2);
[y_closed1, t_closed1] = step(closed_sys1, simtime);
plot(t_closed1, y_closed1, 'b', 'LineWidth', 1.5);
xlabel('t, s');
ylabel('y');
grid on;
title('Замкнутая система 1', 'FontWeight', 'normal');
xlim([0 simtime]);

exportgraphics(gcf, "images/step_combined1.png", Resolution=200);

figure('Position', [100 100 800 600]);
set(gcf, 'Color', 'w');

w = logspace(-10, 10, 10000);
[re1, im1] = nyquist(open_sys1, w);
re1 = squeeze(re1);
im1 = squeeze(im1);

plot(re1, im1, 'b', 'LineWidth', 1.5);
hold on;
plot(re1, -im1, 'g', 'LineWidth', 1.5, 'LineStyle', '-');

arrow_spacing = 4000;
arrow_size = 0.1;

for k = 1:arrow_spacing:length(re1)-1
    x1 = re1(k);
    y1 = im1(k);
    x2 = re1(k+1);
    y2 = im1(k+1);

    theta = atan2(y2 - y1, x2 - x1);

    arrow_x = [0 -0.5 -0.5 0]*arrow_size;
    arrow_y = [0 0.2 -0.2 0]*arrow_size;

    R = [cos(theta) -sin(theta); sin(theta) cos(theta)];
    coords = R * [arrow_x; arrow_y];
    patch(x1 + coords(1,:), y1 + coords(2,:), 'r', 'EdgeColor','none');
end

plot(-1, 0, 'xr', 'MarkerSize', 9, 'LineWidth', 2);
xlabel('Re(W(j\omega))');
ylabel('Im(W(j\omega))');
grid on;
title('Годограф Найквиста 1', 'FontWeight', 'normal');
axis equal;
box on;
ax = gca;
ax.XGrid = 'on';
ax.YGrid = 'on';
ax.GridLineStyle = '-';
ax.GridAlpha = 0.3;
ax.LineWidth = 1.5;

exportgraphics(gcf, "images/nyquist1.png", Resolution=200);

figure('Position', [100 100 1000 800]);
set(gcf, 'Color', 'w');

[mag1, phase1, w1] = bode(open_sys1);
mag1 = squeeze(mag1);
phase1 = squeeze(phase1);

subplot(2,1,1);
semilogx(w1, 20*log10(mag1), 'b', 'LineWidth', 1.5);
grid on;
ylabel('L, dB');
title('ЛАЧХ разомкнутой системы 1', 'FontWeight', 'normal');

subplot(2,1,2);
semilogx(w1, phase1, 'r', 'LineWidth', 1.5);
grid on;
xlabel('\omega, rad/s');
ylabel('\phi, deg');
title('ЛФЧХ разомкнутой системы 1', 'FontWeight', 'normal');

exportgraphics(gcf, "images/bode1.png", Resolution=200);

% Система 2
R2 = [-9 -31 -73 121 432];
Q2 = [1 7 21 33 28 10];
D2 = [1 -2 -10 -40 149 442];
open_sys2 = tf(R2, Q2);
closed_sys2 = tf(R2, D2);
p2 = pole(open_sys2);
z2 = zero(open_sys2);
p_closed2 = pole(closed_sys2);
z_closed2 = zero(closed_sys2);


figure('Position', [100 100 1200 500]);
set(gcf, 'Color', 'w');

subplot(1,2,1);
hold on;
grid on;
xlabel('Re');
ylabel('Im');
title('Разомкнутая система 2', 'FontWeight', 'normal');

for i = 1:length(p2)
    if real(p2(i)) < 0
        plot(real(p2(i)), imag(p2(i)), 'xg', 'MarkerSize', 10, 'LineWidth', 2);
    else
        plot(real(p2(i)), imag(p2(i)), 'xr', 'MarkerSize', 10, 'LineWidth', 2);
    end
end
for i = 1:length(z2)
    plot(real(z2(i)), imag(z2(i)), 'bo', 'MarkerSize', 8, 'LineWidth', 2);
end
xlim([-3 3]);
ylim([-4 4]);
box on;
h1 = plot(NaN, NaN, 'xg', 'MarkerSize', 10, 'LineWidth', 2);
h2 = plot(NaN, NaN, 'xr', 'MarkerSize', 10, 'LineWidth', 2);
h3 = plot(NaN, NaN, 'bo', 'MarkerSize', 8, 'LineWidth', 2);
legend([h1, h2, h3], {'Уст. полюс', 'Неуст. полюс', 'Ноль'}, 'Location', 'best');

subplot(1,2,2);
hold on;
grid on;
xlabel('Re');
ylabel('Im');
title('Замкнутая система 2', 'FontWeight', 'normal');

for i = 1:length(p_closed2)
    if real(p_closed2(i)) < 0
        plot(real(p_closed2(i)), imag(p_closed2(i)), 'xg', 'MarkerSize', 10, 'LineWidth', 2);
    else
        plot(real(p_closed2(i)), imag(p_closed2(i)), 'xr', 'MarkerSize', 10, 'LineWidth', 2);
    end
end
for i = 1:length(z_closed2)
    plot(real(z_closed2(i)), imag(z_closed2(i)), 'bo', 'MarkerSize', 8, 'LineWidth', 2);
end
xlim([-3 5]);
ylim([-4 4]);
box on;
h1 = plot(NaN, NaN, 'xg', 'MarkerSize', 10, 'LineWidth', 2);
h2 = plot(NaN, NaN, 'xr', 'MarkerSize', 10, 'LineWidth', 2);
h3 = plot(NaN, NaN, 'bo', 'MarkerSize', 8, 'LineWidth', 2);
legend([h1, h2, h3], {'Уст. полюс', 'Неуст. полюс', 'Ноль'}, 'Location', 'best');

exportgraphics(gcf, "images/pz_combined2.png", Resolution=200);

figure('Position', [100 100 800 800]);
set(gcf, 'Color', 'w');

simtime = 20;
subplot(2,1,1);
[y_open2, t2] = step(open_sys2, simtime);
plot(t2, y_open2, 'b', 'LineWidth', 1.5);
xlabel('t, s');
ylabel('y');
grid on;
title('Разомкнутая система 2', 'FontWeight', 'normal');
xlim([0 simtime]);

subplot(2,1,2);
[y_closed2, t_closed2] = step(closed_sys2, simtime);
plot(t_closed2, y_closed2, 'b', 'LineWidth', 0.85);
xlabel('t, s');
ylabel('y');
grid on;
title('Замкнутая система 2', 'FontWeight', 'normal');
xlim([0 simtime]);

exportgraphics(gcf, "images/step_combined2.png", Resolution=200);

figure('Position', [100 100 800 600]);
set(gcf, 'Color', 'w');

w = logspace(-10, 10, 10000);
[re2, im2] = nyquist(open_sys2, w);
re2 = squeeze(re2);
im2 = squeeze(im2);

plot(re2, im2, 'b', 'LineWidth', 1.5);
hold on;
plot(re2, -im2, 'g', 'LineWidth', 1.5, 'LineStyle', '-');

arrow_spacing = 5000;
arrow_size = 5;

for k = 1:arrow_spacing:length(re2)-1
    x1 = re2(k);
    y1 = im2(k);
    x2 = re2(k+1);
    y2 = im2(k+1);

    theta = atan2(y2 - y1, x2 - x1);

    arrow_x = [0 -0.5 -0.5 0]*arrow_size;
    arrow_y = [0 0.2 -0.2 0]*arrow_size;

    R = [cos(theta) -sin(theta); sin(theta) cos(theta)];
    coords = R * [arrow_x; arrow_y];
    patch(x1 + coords(1,:), y1 + coords(2,:), 'r', 'EdgeColor','none');
end

plot(-1, 0, 'xr', 'MarkerSize', 9, 'LineWidth', 2);
xlabel('Re(W(j\omega))');
ylabel('Im(W(j\omega))');
grid on;
title('Годограф Найквиста 2', 'FontWeight', 'normal');
axis equal;
box on;
ax = gca;
ax.XGrid = 'on';
ax.YGrid = 'on';
ax.GridLineStyle = '-';
ax.GridAlpha = 0.3;
ax.LineWidth = 1.5;

exportgraphics(gcf, "images/nyquist2.png", Resolution=200);

figure('Position', [100 100 1000 800]);
set(gcf, 'Color', 'w');

[mag2, phase2, w2] = bode(open_sys2);
mag2 = squeeze(mag2);
phase2 = squeeze(phase2);

subplot(2,1,1);
semilogx(w2, 20*log10(mag2), 'b', 'LineWidth', 1.5);
grid on;
ylabel('L, dB');
title('ЛАЧХ разомкнутой системы 2', 'FontWeight', 'normal');

subplot(2,1,2);
semilogx(w2, phase2, 'r', 'LineWidth', 1.5);
grid on;
xlabel('\omega, rad/s');
ylabel('\phi, deg');
title('ЛФЧХ разомкнутой системы 2', 'FontWeight', 'normal');

exportgraphics(gcf, "images/bode2.png", Resolution=200);

% Система 3
R3 = [17 65 271 528 390];
Q3 = [1 -5 9 -3 -8 10];
D3 = [1 12 74 268 520 400];
open_sys3 = tf(R3, Q3);
closed_sys3 = tf(R3, D3);
p3 = pole(open_sys3);
z3 = zero(open_sys3);
p_closed3 = pole(closed_sys3);
z_closed3 = zero(closed_sys3);



figure('Position', [100 100 1200 500]);
set(gcf, 'Color', 'w');

subplot(1,2,1);
hold on;
grid on;
xlabel('Re');
ylabel('Im');
title('Разомкнутая система 3', 'FontWeight', 'normal');

for i = 1:length(p3)
    if real(p3(i)) < 0
        plot(real(p3(i)), imag(p3(i)), 'xg', 'MarkerSize', 10, 'LineWidth', 2);
    else
        plot(real(p3(i)), imag(p3(i)), 'xr', 'MarkerSize', 10, 'LineWidth', 2);
    end
end
for i = 1:length(z3)
    plot(real(z3(i)), imag(z3(i)), 'bo', 'MarkerSize', 8, 'LineWidth', 2);
end
xlim([-2 3]);
ylim([-4 4]);
box on;
h1 = plot(NaN, NaN, 'xg', 'MarkerSize', 10, 'LineWidth', 2);
h2 = plot(NaN, NaN, 'xr', 'MarkerSize', 10, 'LineWidth', 2);
h3 = plot(NaN, NaN, 'bo', 'MarkerSize', 8, 'LineWidth', 2);
legend([h1, h2, h3], {'Уст. полюс', 'Неуст. полюс', 'Ноль'}, 'Location', 'best');

subplot(1,2,2);
hold on;
grid on;
xlabel('Re');
ylabel('Im');
title('Замкнутая система 3', 'FontWeight', 'normal');

for i = 1:length(p_closed3)
    if real(p_closed3(i)) < 0
        plot(real(p_closed3(i)), imag(p_closed3(i)), 'xg', 'MarkerSize', 10, 'LineWidth', 2);
    else
        plot(real(p_closed3(i)), imag(p_closed3(i)), 'xr', 'MarkerSize', 10, 'LineWidth', 2);
    end
end
for i = 1:length(z_closed3)
    plot(real(z_closed3(i)), imag(z_closed3(i)), 'bo', 'MarkerSize', 8, 'LineWidth', 2);
end
xlim([-4 0]);
ylim([-5 5]);
box on;
h1 = plot(NaN, NaN, 'xg', 'MarkerSize', 10, 'LineWidth', 2);
h2 = plot(NaN, NaN, 'xr', 'MarkerSize', 10, 'LineWidth', 2);
h3 = plot(NaN, NaN, 'bo', 'MarkerSize', 8, 'LineWidth', 2);
legend([h1, h2, h3], {'Уст. полюс', 'Неуст. полюс', 'Ноль'}, 'Location', 'best');

exportgraphics(gcf, "images/pz_combined3.png", Resolution=200);

figure('Position', [100 100 800 800]);
set(gcf, 'Color', 'w');

simtime = 10;
subplot(2,1,1);
[y_open3, t3] = step(open_sys3, simtime);
plot(t3, y_open3, 'b', 'LineWidth', 1.5);
xlabel('t, s');
ylabel('y');
grid on;
title('Разомкнутая система 3', 'FontWeight', 'normal');
xlim([0 simtime]);

subplot(2,1,2);
[y_closed3, t_closed3] = step(closed_sys3, simtime);
plot(t_closed3, y_closed3, 'b', 'LineWidth', 1.5);
xlabel('t, s');
ylabel('y');
grid on;
title('Замкнутая система 3', 'FontWeight', 'normal');
xlim([0 simtime]);

exportgraphics(gcf, "images/step_combined3.png", Resolution=200);

figure('Position', [100 100 800 600]);
set(gcf, 'Color', 'w');

w = logspace(-10, 10, 10000);
[re3, im3] = nyquist(open_sys3, w);
re3 = squeeze(re3);
im3 = squeeze(im3);

plot(re3, im3, 'b', 'LineWidth', 1.5);
hold on;
plot(re3, -im3, 'g', 'LineWidth', 1.5, 'LineStyle', '-');

arrow_spacing = 5000;
arrow_size = 5;

for k = 1:arrow_spacing:length(re3)-1
    x1 = re3(k);
    y1 = im3(k);
    x2 = re3(k+1);
    y2 = im3(k+1);

    theta = atan2(y2 - y1, x2 - x1);

    arrow_x = [0 -0.5 -0.5 0]*arrow_size;
    arrow_y = [0 0.2 -0.2 0]*arrow_size;

    R = [cos(theta) -sin(theta); sin(theta) cos(theta)];
    coords = R * [arrow_x; arrow_y];
    patch(x1 + coords(1,:), y1 + coords(2,:), 'r', 'EdgeColor','none');
end

plot(-1, 0, 'xr', 'MarkerSize', 9, 'LineWidth', 2);
xlabel('Re(W(j\omega))');
ylabel('Im(W(j\omega))');
grid on;
title('Годограф Найквиста 3', 'FontWeight', 'normal');
axis equal;
box on;
ax = gca;
ax.XGrid = 'on';
ax.YGrid = 'on';
ax.GridLineStyle = '-';
ax.GridAlpha = 0.3;
ax.LineWidth = 1.5;

exportgraphics(gcf, "images/nyquist3.png", Resolution=200);

figure('Position', [100 100 1000 800]);
set(gcf, 'Color', 'w');

[mag3, phase3, w3] = bode(open_sys3);
mag3 = squeeze(mag3);
phase3 = squeeze(phase3);

subplot(2,1,1);
semilogx(w3, 20*log10(mag3), 'b', 'LineWidth', 1.5);
grid on;
ylabel('L, dB');
title('ЛАЧХ разомкнутой системы 3', 'FontWeight', 'normal');

subplot(2,1,2);
semilogx(w3, phase3, 'r', 'LineWidth', 1.5);
grid on;
xlabel('\omega, rad/s');
ylabel('\phi, deg');
title('ЛФЧХ разомкнутой системы 3', 'FontWeight', 'normal');

exportgraphics(gcf, "images/bode3.png", Resolution=200);
\end{matlabcode}

\matlabtitle{Задание 2}
\begin{matlabcode}

\end{matlabcode}

\matlabtitle{Задание 3}
\begin{matlabcode}

\end{matlabcode}
